\documentclass[a4paper,10pt]{article}
\usepackage[utf8]{inputenc}
\usepackage[T1]{fontenc}
\usepackage[french]{babel}
\usepackage{enumitem}
% Puces pleines au premier niveau, vides au second : jamais de tiret, confondu avec un signe moins.
\frenchsetup{StandardItemLabels=true}
\setlist[itemize,1]{label=\textbullet}
\setlist[itemize,2]{label=$\circ$}
\usepackage[table]{xcolor}
\usepackage{geometry}
\usepackage{tabularx}
\usepackage{xltabular}
\usepackage{amssymb}
\usepackage{amsmath}
\usepackage{textcomp}
\usepackage{fontawesome5}
\usepackage{tikz}

\geometry{hmargin=0.8cm, vmargin=1.5cm}

% Commande 6 cases d'auto-évaluation sur 2 lignes
\newcommand{\caseseval}{%
    \centering
    $\square\;\square\;\square$ \\
    \vspace{2mm}
    $\square\;\square\;\square$
}
\renewcommand{\tabularxcolumn}[1]{m{#1}}

\begin{document}

\begin{center}
    \LARGE \textbf{\textcolor{orange!80!black}{\faBookOpen} \quad FICHE MÉMO}\\[0.5em]
    \large \textbf{{{NIVEAU}} - {{CHAPITRE}} - {{TITRE}}}
\end{center}

\vspace{0.2cm}

% arraystretch : 2.2 par défaut ; 1.8 si la fiche gagne à être plus dense.
% Le tableau est continu : les objectifs s'enchainent et la coupure de page est automatique.
\renewcommand{\arraystretch}{2.2}
\begin{center}
    \begin{xltabular}{\dimexpr\textwidth-5pt\relax}{|m{2cm}|X|X|X|}
        \hline
        \rowcolor{gray!10}
        \centering \scriptsize \textbf{\textcolor{red}{\faTimes}/\textcolor{green!60!black}{\faCheck} Auto-évaluation} &
        \centering \textbf{\textcolor{blue!70!black}{\faQuestionCircle} Question} &
        \centering \textbf{\textcolor{red!70!black}{\faBullseye} Réponse} &
        \centering \textbf{\textcolor{orange!90!black}{\faLightbulb} Stratégie / Conseil} \arraybackslash \\
        \hline
        \endhead

        % --- BLOC OBJECTIF (à répéter pour chaque objectif du cours) ---
        \multicolumn{4}{|c|}{\cellcolor{blue!10}\parbox{\dimexpr\textwidth-20pt\relax}{\centering \textit{\textbf{Objectif : {{INTITULÉ_DE_L_OBJECTIF}}}}}} \\
        \hline

        % --- LIGNE (à répéter pour chaque question atomique) ---
        \caseseval & {{QUESTION_ATOMIQUE_OU_TIKZ}} & {{REPONSE_COURTE}} & \footnotesize \textit{{{FEEDBACK_PROCESSUS}}} \\ \hline

    \end{xltabular}
\end{center}
\end{document}
